\chapter*{如何使用这份讲义}
\addcontentsline{toc}{chapter}{如何使用这份讲义}
这份讲义只讲一件事：在一个有序数组里，怎样用比较次数尽量少的办法找到目标，以及为什么写对它比想象中难。全文只有一章，建议按顺序读完，再回头做章末的自测。

\textbf{全书脉络。}先说清二分查找成立的前提和它为什么快，再把朴素写法一步步改成不会出边界错误的写法，最后把``找一个值''推广为``找第一个满足条件的位置''。后面的内容都建立在这一推广之上。

\textbf{色块含义。}正文里有四种色块：
\begin{keypoint}值得记住的结论或判断方法。\end{keypoint}
\begin{pitfall}手算或编程时最容易错的地方。\end{pitfall}
\begin{errata}原始教材或课件中的错误及其修正。\end{errata}
\begin{supp}超出主线、但有助于理解的补充内容。\end{supp}

\textbf{符号。}专用记号在每节第一次出现时，右侧页边有简短旁注；完整含义见附录的符号表。
